\section{TLA+ спецификация алгоритма консенсуса Raft}

Для формального описания и последующей проверки корректности алгоритма
консенсуса в работе использована спецификация на языке TLA+. Разработанная
спецификация моделирует основные элементы поведения кластера: локальные
состояния узлов, обмен сообщениями между ними, накопление и согласование
записей журнала, а также продвижение подтверждённых записей в машину состояний.
При этом в спецификации сознательно опущены низкоуровневые детали программной
реализации, такие как потоки исполнения, особенности gRPC, структура классов
или формат хранения данных в памяти. Вместо этого фиксируется именно та
семантика, которая существенна для проверки корректности алгоритма на уровне
распределённого взаимодействия.

Следует подчеркнуть, что спецификация используется в работе в двух ролях.
С одной стороны, она выступает как самостоятельная формальная модель Raft,
позволяющая проверить базовые свойства безопасности с помощью TLC. С другой
стороны, она служит эталоном для сопоставления с реальной реализацией:
именно в терминах этой модели впоследствии интерпретируются трассы исполнения
инструментированной программы. Поэтому важно, чтобы спецификация была
достаточно абстрактной для удобной проверки, но при этом достаточно точной,
чтобы по ней можно было судить о корректности поведения реализованного кода.

Полный текст спецификации приведён в приложении Б. В данном разделе
рассматриваются её структура, основные переходы системы, проверяемые свойства и
порядок модельной проверки.

\subsection{Структура спецификации}

Модуль \texttt{Raft} расширяет стандартные модули \texttt{Naturals},
\texttt{FiniteSets}, \texttt{Sequences} и \texttt{TLC}. В спецификации задаются
константы, определяющие множество серверов \texttt{Server}, множество возможных
значений \texttt{Value}, роли узлов (\texttt{Follower}, \texttt{Candidate},
\texttt{Leader}), типы сообщений протокола
(\texttt{RequestVoteRequest}, \texttt{RequestVoteResponse},
\texttt{AppendEntriesRequest}, \texttt{AppendEntriesResponse}), а также
ограничения на число выборов и рестартов
(\texttt{MaxElections}, \texttt{MaxRestarts}). Эти параметры фиксируют модель
кластера и позволяют ограничить пространство состояний при проверке.

Состояние системы описывается набором глобальных и локальных переменных.
Глобальная переменная \texttt{messages} моделирует сеть как мультимножество
сообщений, находящихся ``в полёте''. Дополнительно введены вспомогательные
переменные \texttt{acked}, \texttt{electionCtr} и \texttt{restartCtr}, которые
используются для проверки свойств модели и ограничения числа переходов. Для
каждого сервера задаются стандартные переменные Raft:
\texttt{currentTerm}, \texttt{state}, \texttt{votedFor}, журнал
\texttt{log}, индекс коммита \texttt{commitIndex}, множество уже полученных
голосов \texttt{votesGranted}, а также лидерские структуры \texttt{nextIndex},
\texttt{matchIndex} и \texttt{pendingResponse}. Последняя переменная позволяет
запретить одновременную отправку нескольких \texttt{AppendEntries} одному и тому
же узлу и тем самым уменьшить число межсостояний в модели.

Начальное состояние задаётся оператором \texttt{Init}. В нём все узлы находятся
в состоянии \texttt{Follower}, имеют начальный терм, пустой журнал, нулевой
\texttt{commitIndex}, не голосовали ни за одного кандидата и не имеют
недоставленных сообщений. Таким образом, \texttt{Init} описывает корректную
исходную конфигурацию кластера до начала выборов и репликации.

\subsection{Модель переходов}

Динамика протокола задаётся оператором \texttt{Next}, представляющим собой
дизъюнкцию всех допустимых действий. В модели отражены основные переходы Raft.

Действие \texttt{RequestVote(i)} описывает начало выборов: сервер $i$
переходит в состояние кандидата, увеличивает терм, голосует за себя и рассылает
запросы \texttt{RequestVoteRequest} другим узлам. Если кандидат получает
большинство голосов, срабатывает \texttt{BecomeLeader(i)}: узел становится
лидером и инициализирует структуры \texttt{nextIndex} и \texttt{matchIndex} для
репликации журнала.

Репликация журнала задаётся действием \texttt{AppendEntries(i, j)}, в котором
лидер $i$ отправляет последователю $j$ сообщение \texttt{AppendEntriesRequest}.
В модели за один шаг передаётся не более одной записи журнала, что упрощает
проверку, но сохраняет основную семантику алгоритма. Ответы на такие запросы
обрабатываются переходом \texttt{HandleAppendEntriesResponse}, который либо
обновляет \texttt{matchIndex} и \texttt{nextIndex} при успехе, либо уменьшает
\texttt{nextIndex} при обнаружении расхождения журналов.

Получение и обработка голосов задаются действиями
\texttt{HandleRequestVoteRequest} и \texttt{HandleRequestVoteResponse}. В них
формализуются стандартные правила Raft: голос выдаётся только если кандидат не
уступает получателю по ``свежести'' журнала и если узел ещё не голосовал в
данном терме.

Отдельное действие \texttt{UpdateTerm} отвечает за переход узла в новый терм при
получении сообщения с большим значением \texttt{term}. В этом случае узел
становится \texttt{Follower} и сбрасывает локальное решение о голосовании.
Таким образом, в спецификации явно выделен механизм, обеспечивающий монотонный
рост термов и корректную реакцию на появление более ``свежего'' состояния
кластера.

Продвижение подтверждённого индекса описывает действие
\texttt{AdvanceCommitIndex}. Лидер вычисляет максимальный индекс, который
реплицирован на кворум узлов и относится к его текущему терму, после чего
перемещает \texttt{commitIndex} вперёд. Это соответствует ключевому правилу
Raft, запрещающему коммит записей старых термов только на основании большинства
подтверждений.

Наконец, в модели присутствует переход \texttt{Restart(i)}, который имитирует
перезапуск узла. При рестарте сбрасываются временные структуры состояния
(например, лидерские и кандидатские переменные), однако сохраняются
персистентные данные, прежде всего \texttt{currentTerm}, \texttt{votedFor} и
\texttt{log}. Такое поведение соответствует модели аварийного завершения и
восстановления, заложенной в алгоритме Raft.

\subsection{Проверяемые свойства}

На уровне модели проверяются прежде всего свойства безопасности. В
спецификации заданы два основных инварианта.

Инвариант \texttt{NoLogDivergence} утверждает, что если две реплики считают
некоторую префиксную часть журнала подтверждённой, то эти префиксы совпадают.
Иначе говоря, у двух узлов не может быть разных значений в одной и той же уже
зафиксированной позиции журнала. Это выражает фундаментальное свойство
согласованности реплицированного лога.

Инвариант \texttt{LeaderHasAllAckedValues} формализует требование, согласно
которому любое значение, уже подтверждённое клиенту, должно присутствовать в
журнале любого актуального лидера. Тем самым исключается потеря данных,
которые были признаны системой завершёнными.

Помимо инвариантов, спецификация определяет полную систему переходов
\texttt{Spec == Init /\ [][Next]\_vars}. Для исключения пустых шагов также
используется условие слабой справедливости \texttt{WF\_vars(Next)}, что
позволяет анализировать поведение системы без искусственного застревания в
неизменных состояниях.

\subsection{Проверка модели}

Проверка спецификации выполнялась с помощью средства \texttt{TLC}. В
конфигурации модели задавались конкретные множества серверов и значений, а
также ограничения \texttt{MaxElections} и \texttt{MaxRestarts}, позволяющие
сделать перебор конечным. В ходе проверки анализировались все достижимые
состояния в рамках выбранных параметров, после чего автоматически проверялись
инварианты безопасности.

Таким образом, TLA+-спецификация позволила формально зафиксировать поведение
алгоритма Raft на уровне переходной системы и проверить, что в рассматриваемой
модели не нарушаются ключевые свойства согласованности журнала и сохранности
подтверждённых данных.
